\chapter*{Sources and evidence notes}
\addcontentsline{toc}{chapter}{Sources and evidence notes}
Access checked 24 September 2026. The Caltech record and submitted sticker
manuscript's printed pages 1--2 (PDF pages 2--3), covering representation
and formal operations, were visually inspected. Text extraction from this
older PDF was unreliable. Its physical-implementation sections were not
audited; no selective-clear protocol is claimed here.
The 2026 paper's main-text renewal and scale-up sections were inspected.
The register algebra, AND reference, error calculations and exercises are
original derivations under explicitly idealized operations.
\begingroup\let\chapter\section
\begin{thebibliography}{9}
\bibitem{stickers} S. Roweis, E. Winfree, R. Burgoyne, N. V. Chelyapov,
M. F. Goodman, P. W. K. Rothemund and L. M. Adleman.
A sticker-based model for DNA computation.
\emph{Journal of Computational Biology} 5(4), 615--629 (1998).
DOI: 10.1089/cmb.1998.5.615.
\url{https://authors.library.caltech.edu/records/jvmmc-dtb97}.
Submitted manuscript:
\url{https://authors.library.caltech.edu/records/jvmmc-dtb97/files/stickers_1_.pdf}.
\bibitem{sdc} T. Stérin, A. Eshra, C. G. Evans, J. Adio and D. Woods.
A thermodynamically favoured molecular computer.
\emph{Nature} 657, 646--652 (2026).
\url{https://doi.org/10.1038/s41586-026-10996-5}.
\end{thebibliography}
\endgroup
Assigned retention probabilities are illustrative, not fitted or measured
sticker-system performance. Logical operations do not establish laboratory
feasibility.
